\documentclass[journal]{IEEEtran}
\usepackage{amsmath,amssymb,amsfonts,graphicx,booktabs,array,cite,url}
\usepackage[hidelinks]{hyperref}
\hypersetup{pdftitle={Channel Estimation for Rydberg Receivers Using Learned Hankel Filtering},pdfauthor={Abdullah Haatim},pdfsubject={Simulation study of magnitude-only channel estimation}}
\graphicspath{{figures/}}
\interdisplaylinepenalty=2500
\newcommand{\CN}{\mathcal{CN}}
\newcommand{\Hank}{\mathcal{H}}
\DeclareMathOperator{\sigmoid}{sigmoid}
\DeclareMathOperator{\softplus}{softplus}
\DeclareMathOperator{\tr}{tr}
\hyphenation{Ryd-berg mea-sure-ments}
\begin{document}
\title{Channel Estimation for Rydberg Receivers\\Using Learned Hankel Filtering}
\author{Abdullah~Haatim\thanks{A. Haatim is with the Department of Electrical and Electronics Engineering, Birla Institute of Technology and Science, Pilani, India (e-mail: f20220519@pilani.bits-pilani.ac.in).}}
\maketitle
\begin{abstract}
Rydberg receivers use excited atoms to sense radio signals. Magnitude-only
observations do not directly reveal the wireless channel. Channel structure
can help recover it, but fixed filtering may remove useful details. We
study a small learned estimator that alternates measurement fitting with
Hankel filtering. It learns from withheld observations, without true-channel
targets. It improves on the original rank-search baseline, but a stronger
classical shrinkage method has slightly lower average error and similar
processing time. The learned method makes fewer severe errors in the
tested noise mixture. Weak references and some few-path channels remain
difficult. The result is a compact, label-free alternative with a measured
accuracy--failure trade-off, not a generally superior channel estimator.
\end{abstract}
\begin{IEEEkeywords}
Channel estimation, Hankel filtering, magnitude measurements, Rydberg receivers.
\end{IEEEkeywords}

\section{Introduction}
\IEEEPARstart{R}{ydberg} atoms respond to radio-frequency electric fields,
offering a way to study new wireless receivers. Here, signals from several
users combine with a known reference. The receiver measures their combined
magnitude, not their full complex value. The channel describes changes in
signal strength and phase on the way to the receiver. Model-based methods
and learned methods such as URformer estimate it~\cite{xu2025,xiao2026}.

A few propagation paths can create a repeated pattern across a regular
array. Arranging channel samples in a Hankel matrix lets an estimator use
this pattern. Strong filtering can remove real details; weak filtering can
leave noise. Searching many filtering settings also takes time. We study a
small model that learns how much filtering to apply. The true channel is
never shown to it during training. Instead, some measurements are hidden
from the estimator and used to check its predictions. This is called
self-supervised learning.

Our contributions are: (i) a 188-parameter estimator that combines Rician
measurement updates with smooth Hankel filtering; (ii) training this
estimator on withheld measurements, without channel targets or channel-error
checkpoint selection; and (iii) paired tests of accuracy, processing time
and simpler versions, including fresh-test replication and input-removal
tests. A controlled in-loop OptShrink comparison identifies a trade-off
between average error and severe failures. We report losses alongside gains.

\begin{figure*}[!t]
\centering\includegraphics[width=\textwidth]{fig1_method.pdf}
\caption{Estimator overview. The controller chooses filtering strength and mixing. During training, nine withheld pilots score predictions but do not enter initialization or the ten stages.}
\label{fig:method}
\end{figure*}

\section{Related Work}
Training on withheld measurements is established in image
reconstruction~\cite{yaman2020} and self-supervised denoising is used for
Rydberg sensor signals~\cite{liu2026}. Huang et al.\ learn Hankel-factor
updates from fully sampled targets~\cite{huang2021}. Luo et al.\ combine
self-supervision and structured low-rank reconstruction for complex MRI
data~\cite{luo2023}. DeepTensor fits network-based low-rank factors to an
observed object~\cite{saragadam2022}. Rician-likelihood training for
magnitude data is also established~\cite{parker2023}. Neither
self-supervision, learned Hankel processing nor the Rician loss is new
on its own.

We recover a complex multiuser channel from noisy magnitudes and a known
reference, rather than restoring measured complex samples. A small
controller sets filtering strength instead of generating low-rank factors.
Our companion manuscript studies classical Hankel estimation and rank
selection~\cite{haatimcompanion}; here we learn filtering without channel labels.

Structured alternating least squares~\cite{zachariah2012},
data-driven OptShrink~\cite{nadakuditi2014}, low-rank Hankel denoising
(LRHD)~\cite{yin2021} and robust Hankel completion~\cite{cai2023}
already provide other structural methods. Their original linear or direct
matrix observations differ from our referenced magnitudes. Sparse
amplitude-only beam training is also relevant~\cite{wang2026}, but uses
adaptive sensing without our additive reference. We do not claim an
optimal shrinker or superiority over these complete methods. Our stronger
control puts OptShrink inside Rician EM iterations. This adapts data-driven
shrinkage to the same observations; it is not a reproduction of the full
LRHD algorithm for directly observed Hankel data.

\section{Receiver and Channel Model}
An array with $N$ elements receives $P$ known training symbols, or pilots,
from each of $K$ users. The observations are
\begin{equation}
 Z=|GS+B+W|,\qquad W_{np}\sim\CN(0,v).
 \label{eq:model}
\end{equation}
Here $G\in\mathbb C^{N\times K}$ is the channel,
$S\in\mathbb C^{K\times P}$ contains the pilots, and
$B\in\mathbb C^{N\times P}$ is the known reference. Magnitudes act
entrywise. Independent noise has total complex variance $v$, or $v/2$ in
each real component. The reference makes magnitudes depend on relative
signal phases, but does not ensure accurate recovery in every case.
Reference calibration and noise variance are known, except in stated
input-mismatch checks. Receiver conversion gain is absorbed into the scaling.
This effective model does not simulate atomic dynamics or optical readout.

For the half-wavelength uniform linear array,
\begin{equation}
 g_{nk}=\sum_{\ell=1}^{L_k}\alpha_{\ell k}e^{-j\pi n\sin\theta_{\ell k}},
 \quad n=0,\ldots,N-1.
\end{equation}
Independent gains follow $\CN(0,1/L_k)$, giving each user unit average
channel power. Angles are uniform on $[-\pi/2,\pi/2]$; pilot entries
follow $\CN(0,1)$. The reference has a known spatial response and constant waveform. Define
$\mathrm{SNR}=10\log_{10}(K/v)$ and
$\mathrm{RSR}=10\log_{10}|\alpha_b|^2$, where $\alpha_b$ is reference
amplitude. RSR uses one user's power: for three users, RSR 12~dB is
about 7.23~dB relative to their combined power.

\subsection{What the Measurements Can Determine}
For one array element, write $\mu_p=g^Ts_p+b_p$ and
$x=[\operatorname{Re}g^T,\operatorname{Im}g^T]^T$. If $\mu_p\ne0$, the
real Jacobian of $|\mu_p|$ has row
\begin{equation}
 A_p=\left[\operatorname{Re}c_p^T,-\operatorname{Im}c_p^T\right],
 \qquad c_p=\mu_p^*s_p/|\mu_p|.
\end{equation}
Full column rank of $A$ ensures local recovery; a small minimum singular
value means poor stability. Thus $P\ge2K$ is necessary for this condition,
not sufficient. Full row rank of $S$ alone is insufficient. With $B=0$,
channel rotation leaves magnitudes unchanged. A known reference can break
this ambiguity, but its strength alone does not ensure recovery. This
unrestricted-row condition is local, not a global or sparse-channel
guarantee~\cite{gao2018}. Proof and information matrix are in the supplement.

\section{Proposed Estimator}
The estimator starts with a rough channel guess. It repeats two steps ten
times: one pulls the guess toward the measurements, and the other smooths
it using the array structure. A small network decides how much smoothing
to apply at each step (Fig.~\ref{fig:method}).

\subsection{Measurement Update}
The magnitude of a complex signal with Gaussian noise follows a Rician
distribution. For noiseless magnitude $a=|\mu|$, its density is
\begin{equation}
 p(z\mid a,v)=\frac{2z}{v}e^{-(z^2+a^2)/v}I_0(2za/v),
\end{equation}
where $I_0$ is a modified Bessel function. Let $M=GS+B$ and
$\kappa=2Z|M|/v$. Entrywise, form
\begin{equation}
 Q=Z\frac{I_1(\kappa)}{I_0(\kappa)}\frac{M}{|M|},
\end{equation}
using the continuous limit at $M=0$. The expectation-maximization (EM)
update and learned step are
\begin{align}
 G_{\rm EM}&=(Q-B)S^H(SS^H)^{-1},\\
 G_{\rm dc}&=G+\gamma_t(G_{\rm EM}-G),\quad
 \gamma_t=2\sigmoid(a_t).
\end{align}
A linear solve evaluates the inverse. Input pilots must have full row
rank. Ten stages are a fixed computation budget, not a proved optimum.

\subsection{Smooth Hankel Filtering}
Each updated user channel forms a Hankel matrix $H=\Hank(g_{{\rm dc},k})$,
with equal entries along each antidiagonal. At $N=32$, it is $17\times16$.
Singular values describe component strengths. Hard rank truncation deletes
the smaller components. Our smooth rule is
\begin{equation}
 \widehat\sigma_i=\frac{\sigma_i^3}{\sigma_i^2+\tau^2}.
\end{equation}
The multiplier $\sigma_i^2/(\sigma_i^2+\tau^2)$ is near one for large
components and near zero for small ones. This can reduce noise but also
bias genuine weak paths. Smoothness and a stable matrix solve motivate
the choice; they do not make it an optimal denoiser.

Set $q=\lceil N/2\rceil$, $s_H=\|H\|_F/\sqrt q$, and
$\tau=s_H[0.001+\softplus(u)]$. The equivalent operation is
\begin{align}
 \widehat H&=H-\tau^2H(H^HH+\tau^2I)^{-1},\\
 g_k^{t+1}&=g_{{\rm dc},k}+\eta_{t,k}
 (\Hank^\dagger\widehat H_k-g_{{\rm dc},k}),
\end{align}
where $\Hank^\dagger$ averages antidiagonals and
$\eta_{t,k}=\sigmoid(w)$. The code normalizes by $s_H$ before solving and
guards zero scales. Gradients pass through the solve, without an
identity-backward approximation. Neither control estimates path count.
The filter does not increase the Hankel Frobenius norm, but this does not imply
smaller channel error. The relaxed EM step decreases its quadratic
surrogate for $0\le\gamma_t\le2$; the following filter need not decrease
the likelihood. We make no convergence claim for the full estimator.

\subsection{Controller Inputs}
The network has eight inputs, 16 hidden units, two outputs and SiLU
activation $x\sigmoid(x)$. Its 178 parameters and ten step sizes give 188. For
$p_i=\sigma_i/\sum_j\sigma_j$, its first
three inputs are
\begin{equation}
 e^{-\sum_i p_i\log p_i}/q,\qquad p_1,\qquad
 \sum_{i=q-\lceil q/2\rceil+1}^{q}p_i.
\end{equation}
They tell the network how spread out the components are, how much the
largest contributes, and how much the smaller components contribute.
Two more inputs measure disagreement with observations and expected
noise relative to estimated channel power:
\begin{align}
 f_4&=\log\left(1+\frac{\operatorname{mean}[(Z-|G_{\rm dc}S+B|)^2]}{v}\right),\\
 f_{5,k}&=\log\left(1+\frac{v\tr[(SS^H)^{-1}]}{K\widehat p_k}\right),
 \quad\widehat p_k=\|g_{{\rm dc},k}\|^2/N.
\end{align}
Inputs $P/(2K)$, $\log q$ and $(t+1)/10$ give measurement count, size and
stage ($t=0,\ldots,9$). Denominator guards avoid division by zero.
True channels, path counts and family labels are not supplied. Weights stay fixed.

\subsection{Training Without True Channels}
Of 30 pilot columns, 21 enter the estimator and nine are withheld.
Training scores predicted withheld magnitudes using the Rician loss
\begin{equation}
 \ell(z,\mu,v)=\frac{(z-|\mu|)^2}{v}-\log I_{0e}(2z|\mu|/v).
\end{equation}
Here $I_{0e}$ is the scaled Bessel function. We omit an observation-only
term and use small-field series at continuous limits. The loss fits
Rician amplitude, not its mean. Splits change each epoch; separate
validation cases with a fixed split select the checkpoint. Deployment
uses all pilots, without extra measurements or weight training.
This avoids test-driven selection; it does not make withheld likelihood
equivalent to channel NMSE or remove dependence on the validation mixture.

\section{Simulation Study}
\subsection{Setup and Comparators}
Table~\ref{tab:setup} gives the settings. Half the training cases have
three to seven paths per user; the rest use fixed counts of 2, 4, 8, 12
and 16. Validation weights 126 SNR, RSR and path-setting combinations
equally. These condition labels set validation weights only; model
selection uses withheld likelihood, not channel error.

\begin{table}[!t]
\caption{Main Simulation and Training Settings}\label{tab:setup}
\centering\small
\begin{tabular}{@{}ll@{}}\toprule
Setting & Value\\\midrule
Array elements / users / pilots & $32/3/30$\\
Training / validation cases & $16{,}384/2{,}048$\\
SNR values (dB) & $-10,-5,0,5,10,15,20$\\
Training RSR values (dB) & $5,12,20$\\
Training seeds & 3, sharing training cases\\
Estimator stages / parameters & $10/188$\\
Training epochs / batch size & $20/32$\\
Optimizer / learning rate & Adam / $0.001$\\
Input / withheld pilots & $21/9$\\
Training / evaluation precision & Single / double\\\bottomrule
\end{tabular}
\end{table}

The main study trains 18 models: two updates, three filtering choices and
three seeds. Seeds change weights, batch order and pilot splits, but share
training cases. Training took 3.67 hours with overlapping workers,
excluding data preparation and preflight. Tests cover 32 conditions with
400 paired cases each. Two noise-input tests reuse cases, giving 12,000
distinct cases. All methods share channels, pilots and noise per case.

HS-GS, the companion study's classical Hankel estimator~\cite{haatimcompanion},
fits ranks $1,\ldots,q$ on 70\% of pilots, using 25 EM updates and four
rank/Hankel projection sweeps per update. Squared magnitude residuals on
the remaining pilots select a rank. It then refits all pilots for 100
updates. Rank search uses no extra observations and counts toward time.
Its rank-selection loss differs from our Rician training loss.
Other controls use EM, Rician or amplitude-gradient updates. Global
filtering has 30 parameters and scales its threshold with $s_H$; no
filtering has ten. Hard-Hankel Rician uses development-selected rank ten
and one projection sweep per update. Constant-output controls replace
adaptive outputs with validation averages without retraining.

The earlier URformer-style experiment uses a different protocol and is
not a matched comparator here. Its scores are not pooled with these results.

Channel error is normalized mean squared error (NMSE):
\begin{equation}
 \mathrm{NMSE}_{\rm dB}=10\log_{10}
 \frac{\sum_i\|\widehat G_i-G_i\|_F^2}{\sum_i\|G_i\|_F^2}.
\end{equation}
Lower is better. Per-seed dB scores and the seven main SNR settings are
averaged equally. Positive gain means lower error. A severe error has
more energy than the true channel (per-case linear NMSE above one).
Pointwise 95\% intervals use 4,000 paired whole-case resamples, conditional
on saved models, excluding training-dataset uncertainty. Classical stopping
methods have a 1,000-update cap; reaching it need not mean convergence.

\begin{figure}[!t]
\centering\includegraphics[width=\columnwidth]{fig2_snr.pdf}
\caption{Main SNR test: $N=32$, $K=3$, $P=30$, RSR 12~dB, and three-to-seven-path channels. Lines average per-seed scores; bands are conditional paired-case 95\% intervals.}
\label{fig:snr}
\end{figure}

\subsection{Accuracy and Processing Time}
\begin{table}[!t]
\caption{Original-Study Summary; Full Comparator Table in Supplement}\label{tab:all}
\centering\footnotesize\setlength{\tabcolsep}{3pt}
\begin{tabular}{@{}lrrr@{}}\toprule
Method & NMSE & Time & Severe\\
 & (dB) & (ms) & (\%)\\\midrule
EM adaptive (primary) & -14.42 & 8.24 & 0.45 \\
EM learned steps only & -10.41 & 2.98 & 22.11 \\
EM global filtering & -12.43 & 5.61 & 19.19 \\
GD adaptive (secondary) & -14.59 & 8.00 & 0.85 \\
HS-GS, rank search & -12.74 & 525.62 & 14.25 \\
EM, 100 updates & -10.28 & 17.01 & 23.57 \\
Rician GD, hard Hankel & -11.06 & 47.39 & 20.93 \\
\bottomrule\end{tabular}
\parbox{\columnwidth}{\footnotesize NMSE and severe-error rates average seven primary SNR conditions equally. Time is the median at SNR 5~dB, RSR 12~dB. Learned scores average three seeds; classical cases are not counted three times.}
\end{table}

HS-GS is the strongest original classical comparator on the declared average
(Table~\ref{tab:all}). The adaptive gradient-descent (GD) core has lower
test NMSE than the primary EM core. We retain EM because validation
likelihood selected it before test generation.

The EM model's average gain over HS-GS is 1.675~dB, with interval
$[1.618,1.730]$~dB and seed-specific gains of 1.533--1.787~dB
(Fig.~\ref{fig:snr}). At 20~dB, the gain is 0.139~dB with interval
$[-0.017,0.294]$~dB, which does not establish a clear improvement.

Severe errors fall from 14.25\% for HS-GS to 0.45\% for the proposed method.

Median time is 8.24~ms versus 525.62~ms for HS-GS, about 64 times faster
(Table~\ref{tab:all}). Timing includes initialization and rank search,
using one Apple M2 CPU thread, double precision, batch size one,
100 cases and three repeats at SNR 5~dB and RSR 12~dB.
These measurements depend on the implementation: PyTorch for learned
methods and NumPy/SciPy for classical methods. Offline training is separate.
With $m=N-q+1$, one adaptive stage costs
$O(NKP+NK^2+K^3+Kmq^2+Kq^3)$: the spectral features still require an SVD.
HS-GS repeats projections across ranks and updates. A small weight count
does not imply low computation.

\begin{figure*}[!t]
\centering\includegraphics[width=\textwidth]{fig3_limits.pdf}
\caption{Limits at SNR 5~dB: (a) reference strength with three-to-seven-path channels; (b) fixed path counts at RSR 12~dB. Positive gain favors the proposed estimator. Bars are conditional pointwise paired-case 95\% intervals. Path counts overlap training.}
\label{fig:limits}
\end{figure*}

\subsection{Operating Limits}
The proposed method loses 1.617~dB at RSR 0~dB, 1.104~dB with two paths,
and 0.203~dB with three paths (Fig.~\ref{fig:limits}). Good average
performance does not ensure good performance with a weak reference or
a simple channel.

Transfer tests at SNR 5~dB and RSR 12~dB give gains for variance inputs
wrong by $\pm3$~dB and for clustered or ten-path alternatives with matched
energy. One-factor dimension tests cover $N\in\{8,16,32,64\}$,
$K\in\{2,3,4\}$ and $P\in\{10,20,30\}$. At $N=8$ the gain interval
includes zero. A new 32-case-per-cell transfer check extends to
$P=6=2K$: at SNR 5~dB, learned NMSE is $-2.197$~dB, compared with
$-14.332$~dB at $P=30$. Full local Jacobian rank does not prevent poor
finite-noise accuracy. The supplement defines all alternatives and reports
dimension results. Path counts overlap training.

\subsection{Further Comparisons}
A separate 15-model follow-up uses two fresh test sets and six adjusted
primary comparisons. Two independent training datasets, one seed each,
give gains of 1.939 and 1.837~dB over HS-GS. Replication is still limited.
A matched supervised recipe is worse on the declared average but better
in some cells; its gap changes with selection and averaging. This does
not establish general superiority of self-supervision. Repeated filtering
gains about 0.70~dB over a retrained final-only model at extra cost.
Mixing stays near one: the controller mainly changes shrinkage strength,
not a reliable on/off decision. All follow-up results remain in the supplement.

\subsection{Which Inputs Help?}
Input-removal controls train on 4,096 shared cases for 12 epochs with
three seeds; two fresh tests total 2,304 worlds. Full inputs improve
average error over spectrum-only, noise/residual-only and size/stage-only
controllers. However, they lose 0.39--0.40~dB against noise/residual-only
at SNR 5~dB, failing the declared 0.25~dB protection rule, and
0.82--0.93~dB on two-path cases. Both input groups help the average,
not every channel. Size/stage inputs, 188 stored parameters, and the
physical update/filter remain in all controls; unused input weights are masked.

\subsection{A Stronger In-Loop Shrinkage Comparison}
A separate controlled study uses 224 development cases and 1,152 fresh
test cases: 128 at each of seven SNRs and two stress conditions. The learned
models remain frozen. OptShrink rank, mixing and iteration count are chosen
on withheld Rician likelihood, separately by known noise level, before test
generation. The same grid includes a final-only control and no filtering.
No true channel enters selection; all methods use the same test measurements.

\begin{table}[!t]
\caption{Fresh stronger-baseline comparison, seven-SNR average}
\label{tab:strong}\centering\small
\begin{tabular}{@{}lrr@{}}\toprule
Method & NMSE (dB) & Severe (\%)\\\midrule
Rician-ranked HS-GS & $-12.890$ & 13.06\\
Final-only OptShrink & $-14.374$ & 3.79\\
In-loop OptShrink & $-14.462$ & 5.02\\
Learned ten-stage model & $-14.319$ & 0.89\\\bottomrule
\end{tabular}
\end{table}

In-loop OptShrink wins the declared average by 0.143~dB, with conditional
paired interval $[0.086,0.198]$~dB (Table~\ref{tab:strong}). The learned model
has fewer severe errors; these occur at SNR $-10$~dB in this sample.
Mean per-case linear NMSE favors the learned model (0.197 versus 0.215),
showing that the ranking depends on the declared metric. This does not
reverse the primary average-dB result. At SNR 5~dB, mean end-to-end times
are 5.67 and 6.20~ms for learned and in-loop methods, including noise-based
setting lookup. Thus the large original speed ratio does not apply to
this stronger baseline. At a weak reference, longer OptShrink runs also
beat the learned model; transferring a stopping rule can fail.

\section{Discussion and Conclusion}
An additional study trains 25 models on 4,096-case datasets for 12 epochs
and tests 3,584 fresh worlds. Five stages lose 1.977~dB against ten;
twenty gain 0.181~dB but take 15.60 rather than 8.04~ms. Width 32 and
two experts gain only 0.065 and 0.023~dB. A strong-to-weak-reference
curriculum gains 0.206~dB on average but loses 2.353~dB in its worst cell.
No variant improves every condition. The supplement records selected/final
scores and independent-data variation; these exploratory tests do not
replace the original model.

An exploratory static noise/reference-amplitude correction helps some
input errors but hurts some weak-reference cases; it cannot determine
unknown reference phase. A nearly noiseless stress test at SNR 60~dB,
far outside training, also exposes filtering bias: at RSR 12~dB the
learned estimate reaches only $-6.647$~dB NMSE, versus $-65.715$~dB for
EM100. Disabling filtering with the same learned steps gives
$-44.020$~dB in a post-hoc diagnostic. Low-noise consistency is not
established.

A small model can learn useful Hankel filtering without true-channel
targets. Its practical setting is repeated estimation near the simulated
training population, where offline training is reused. It is better than
the original rank-search baseline, but not the stronger in-loop control
on average dB error. Its lower severe-error frequency is a separate,
conditional finding, not evidence of general superiority.

All results are simulations. Neither global recovery, full-model
convergence nor superiority over published URformer is established.
Case-resampling intervals exclude training-data uncertainty. The new curriculum
includes weak references but does not remove all failures. Atomic dynamics,
optical readout, unknown reference drift and two-dimensional arrays remain
untested. The supplement retains full comparisons and numerical safeguards.
\bibliographystyle{IEEEtran}
\bibliography{references}
\end{document}
